\section{Introduction}

\lettrine{T}{echnology} scaling beyond 5nm has introduced significant challenges in power delivery and interconnect performance. Although transistor dimensions continue to shrink, the scaling of interconnects has not kept pace, leading to increased wire resistance and tighter voltage margins. Conventional front-side power delivery networks (PDNs), which share routing resources with signal interconnects, face growing difficulty in meeting power integrity requirements without degrading routing efficiency and overall performance\cite{8993617}.

To address such limitations buried power rails (BPR) have been proposed as an early innovation. By relocating power rails beneath active devices and utilizing high-aspect-ratio ruthenium structures, BPR reduces local resistance and frees front-side routing tracks, enabling improved IR-drop characteristics\cite{8993617}. Backside Power Delivery Networks (BS-PDN) move the power grid to the wafer backside which decouples the power and signal routing\cite{10185208}. Experimental findings have shown substantial reductions in voltage droop and measurable frequency improvements in high-performance cells \cite{10185208}. It was also demonstrated significant area savings and enhanced routing efficiency were achieved in both high-performance and high-density scenarios\cite{10631463}, while models indicate more than a four times reduction in IR-drop compared to conventional PDN\cite{8932458}.

These architectural innovations represent a fundamental shift in power delivery design, moving from incremental BEOL optimizations toward structural reorganization of the power distribution network. This literature review therefore examines the evolution from conventional front-side PDNs to buried power rails and backside power delivery, evaluating their signal integrity and power–performance–area benefits as well as the emerging design challenges they introduce in sub-5nm CMOS technologies.